% GENERATED FILE. Edit canonical lessons, then run npm run generate:books.
% canonical-source-hash: fnv1a64:8ab72b28902bf185
% canonical-lessons: AR-C516-tamaddada, AR-C516-inhana, AR-C516-irtajafa, AR-C516-saala, AR-C516-iltiwa

\chapter{Stretching and shaking}
\label{ch:ar-stretching-and-shaking}
\hlchaptermodality{516}

\begin{chapteropening}
\textbf{By the end of this chapter:} \emph{I can say that someone stretched out, bent down, trembled or had a cough, and talk about a sprain.}

Complete the last lesson of chapter 516: I can say that someone stretched out, bent down, trembled or had a cough, and talk about a sprain.
\end{chapteropening}

\section[\texorpdfstring{\ar{تمدد} (tamaddada)}{تمدد (tamaddada)}]{\ar{تمدد} (tamaddada) — to stretch out}

\label{lesson:AR-C516-tamaddada}



\begin{quote}
Before the new one: say the Arabic for to get up, then the Arabic for fainting.
\end{quote}

\subsection*{You'll want to know: \ar{تمدد}}
\textbf{\ar{تمدد}} — \emph{tamaddada} — \textquotedblleft{}to stretch out\textquotedblright{}.

\begin{quote}
\ar{تمدد} \ar{القط} \ar{في} \ar{الشمس}.
\end{quote}

\begin{quote}
\emph{tamaddada al-qiṭṭ fī ash-shams.}
\end{quote}

\begin{quote}
The cat stretched out in the sun.
\end{quote}

Present: \ar{يتمدد}, \emph{yatamaddad}.

\begin{cousinweb}
\textbf{\ar{تمدد}} (\emph{tamaddada}) is built on the native Arabic root \emph{m-d-d}, \textquotedblleft{}to stretch, to extend\textquotedblright{} (Lane). Mudda, a period, which you know, is from the same root: a stretch of time.
\end{cousinweb}

\begin{grammarlens}[title={What you've built}]
One more word for describing what you lived and felt.
\end{grammarlens}

\subsection*{Your turn}
\begin{itemize}
\raggedright
  \item \emph{Say it:} \emph{tamaddada}
  \item \emph{Say it:} \emph{tamaddada}, once more
  \item \emph{Say it:} \emph{tamaddada al-qiṭṭ fī ash-shams.} — The cat stretched out in the sun.
  \item \emph{Recall:} say the Arabic for to get up, then the Arabic for fainting, then say \emph{tamaddada} again
\end{itemize}

\begin{tcolorbox}[breakable,colback=teal!4,colframe=teal!35!black,title={Before you move on}]
What does \ar{تمدد} mean? (\textquotedblleft{}To stretch out\textquotedblright{}.) Say it once more.
\end{tcolorbox}
\section[\texorpdfstring{\ar{انحنى} (inḥanā)}{انحنى (inḥanā)}]{\ar{انحنى} (inḥanā) — to bend down}

\label{lesson:AR-C516-inhana}



\begin{quote}
Before the new one: say the Arabic for fainting, then the Arabic for to stretch out.
\end{quote}

\subsection*{You'll want to know: \ar{انحنى}}
\textbf{\ar{انحنى}} — \emph{inḥanā} — \textquotedblleft{}to bend down\textquotedblright{}.

\begin{quote}
\ar{انحنى} \ar{الرجل} \ar{وأخذ} \ar{المفتاح} \ar{من} \ar{الأرض}.
\end{quote}

\begin{quote}
\emph{inḥanā ar-rajul wa-akhadha al-miftāḥ min al-arḍ.}
\end{quote}

\begin{quote}
The man bent down and took the key from the floor.
\end{quote}

Present: \ar{ينحني}, \emph{yanḥanī}.

\begin{cousinweb}
\textbf{\ar{انحنى}} (\emph{inḥanā}) is built on the native Arabic root \emph{ḥ-n-w}, \textquotedblleft{}to bend, to curve\textquotedblright{} (Lane). Form VII inḥanā is to bend oneself.
\end{cousinweb}

\begin{grammarlens}[title={What you've built}]
One more word for describing what you lived and felt.
\end{grammarlens}

\subsection*{Your turn}
\begin{itemize}
\raggedright
  \item \emph{Say it:} \emph{inḥanā}
  \item \emph{Say it:} \emph{inḥanā}, once more
  \item \emph{Say it:} \emph{inḥanā ar-rajul wa-akhadha al-miftāḥ min al-arḍ.} — The man bent down and took the key from the floor.
  \item \emph{Recall:} say the Arabic for fainting, then the Arabic for to stretch out, then say \emph{inḥanā} again
\end{itemize}

\begin{tcolorbox}[breakable,colback=teal!4,colframe=teal!35!black,title={Before you move on}]
What does \ar{انحنى} mean? (\textquotedblleft{}To bend down\textquotedblright{}.) Say it once more.
\end{tcolorbox}
\section[\texorpdfstring{\ar{ارتجف} (irtajafa)}{ارتجف (irtajafa)}]{\ar{ارتجف} (irtajafa) — to tremble}

\label{lesson:AR-C516-irtajafa}



\begin{quote}
Before the new one: say the Arabic for to stretch out, then the Arabic for to bend down.
\end{quote}

\subsection*{You'll want to know: \ar{ارتجف}}
\textbf{\ar{ارتجف}} — \emph{irtajafa} — \textquotedblleft{}to tremble\textquotedblright{}.

\begin{quote}
\ar{ارتجفت} \ar{يد} \ar{الطالب} \ar{من} \ar{التوتر}.
\end{quote}

\begin{quote}
\emph{irtajafat yad aṭ-ṭālib min at-tawattur.}
\end{quote}

\begin{quote}
The student's hand trembled from nerves.
\end{quote}

Present: \ar{يرتجف}, \emph{yartajif}.

\begin{cousinweb}
\textbf{\ar{ارتجف}} (\emph{irtajafa}) is built on the native Arabic root \emph{r-j-f}, \textquotedblleft{}to quake, to shake\textquotedblright{} (Lane). Form VIII irtajafa is to shake all over.
\end{cousinweb}

\begin{grammarlens}[title={What you've built}]
One more word for describing what you lived and felt.
\end{grammarlens}

\subsection*{Your turn}
\begin{itemize}
\raggedright
  \item \emph{Say it:} \emph{irtajafa}
  \item \emph{Say it:} \emph{irtajafa}, once more
  \item \emph{Say it:} \emph{irtajafat yad aṭ-ṭālib min at-tawattur.} — The student's hand trembled from nerves.
  \item \emph{Recall:} say the Arabic for to stretch out, then the Arabic for to bend down, then say \emph{irtajafa} again
\end{itemize}

\begin{tcolorbox}[breakable,colback=teal!4,colframe=teal!35!black,title={Before you move on}]
What does \ar{ارتجف} mean? (\textquotedblleft{}To tremble\textquotedblright{}.) Say it once more.
\end{tcolorbox}
\section[\texorpdfstring{\ar{سعل} (saʿala)}{سعل (saʿala)}]{\ar{سعل} (saʿala) — to have a cough}

\label{lesson:AR-C516-saala}



\begin{quote}
Before the new one: say the Arabic for to bend down, then the Arabic for to tremble.
\end{quote}

\subsection*{You'll want to know: \ar{سعل}}
\textbf{\ar{سعل}} — \emph{saʿala} — \textquotedblleft{}to have a cough\textquotedblright{}.

\begin{quote}
\ar{سعل} \ar{المريض} \ar{كثيرا} \ar{في} \ar{الليل}.
\end{quote}

\begin{quote}
\emph{saʿala al-marīḍ kathīran fī al-layl.}
\end{quote}

\begin{quote}
The patient coughed a lot in the night.
\end{quote}

Present: \ar{يسعل}, \emph{yasʿul}.

The noun \ar{سعال}, a cough, shares the root.

\begin{cousinweb}
\textbf{\ar{سعل}} (\emph{saʿala}) is built on the native Arabic root \emph{s-ʿ-l}, \textquotedblleft{}to cough\textquotedblright{} (Lane). Suʿāl, a cough, which you know, is its noun.
\end{cousinweb}

\begin{grammarlens}[title={What you've built}]
One more word for describing what you lived and felt.
\end{grammarlens}

\subsection*{Your turn}
\begin{itemize}
\raggedright
  \item \emph{Say it:} \emph{saʿala}
  \item \emph{Say it:} \emph{saʿala}, once more
  \item \emph{Say it:} \emph{saʿala al-marīḍ kathīran fī al-layl.} — The patient coughed a lot in the night.
  \item \emph{Recall:} say the Arabic for to bend down, then the Arabic for to tremble, then say \emph{saʿala} again
\end{itemize}

\begin{tcolorbox}[breakable,colback=teal!4,colframe=teal!35!black,title={Before you move on}]
What does \ar{سعل} mean? (\textquotedblleft{}To have a cough\textquotedblright{}.) Say it once more.
\end{tcolorbox}
\section[\texorpdfstring{\ar{التواء} (iltiwāʾ)}{التواء (iltiwāʾ)}]{\ar{التواء} (iltiwāʾ) — a sprain}

\label{lesson:AR-C516-iltiwa}



\begin{quote}
Before the new one: say the Arabic for to tremble, then the Arabic for to have a cough.
\end{quote}

\subsection*{You'll want to know: \ar{التواء}}
\textbf{\ar{التواء}} — \emph{iltiwāʾ} — \textquotedblleft{}a sprain\textquotedblright{}.

\begin{quote}
\ar{عنده} \ar{التواء} \ar{في} \ar{الكاحل}.
\end{quote}

\begin{quote}
\emph{ʿindahu iltiwāʾ fī al-kāḥil.}
\end{quote}

\begin{quote}
He has a sprained ankle.
\end{quote}

\begin{cousinweb}
\textbf{\ar{التواء}} (\emph{iltiwāʾ}) is built on the native Arabic root \emph{l-w-y}, \textquotedblleft{}to twist, to bend\textquotedblright{} (Lane). Iltiwāʾ is the verbal noun of Form VIII: a twisting.
\end{cousinweb}

\begin{grammarlens}[title={What you've built}]
One more word for describing what you lived and felt.
\end{grammarlens}

\subsection*{Your turn}
\begin{itemize}
\raggedright
  \item \emph{Say it:} \emph{iltiwāʾ}
  \item \emph{Say it:} \emph{iltiwāʾ}, once more
  \item \emph{Say it:} \emph{ʿindahu iltiwāʾ fī al-kāḥil.} — He has a sprained ankle.
  \item \emph{Recall:} say the Arabic for to tremble, then the Arabic for to have a cough, then say \emph{iltiwāʾ} again
\end{itemize}

\begin{tcolorbox}[breakable,colback=teal!4,colframe=teal!35!black,title={Before you move on}]
What does \ar{التواء} mean? (\textquotedblleft{}A sprain\textquotedblright{}.) Say it once more.
\end{tcolorbox}
